\chapter{Code Examples} \label{app:code-examples}
\begin{listing}[p]
    \begin{minted}{c}
const ct_protocol_impl_t quic_protocol_interface =
    {
        .name = "QUIC",
        .protocol_enum = CT_PROTOCOL_QUIC,
        .supports_alpn = true,
        .selection_properties =
        {
            .list =
            {
                [RELIABILITY] = {.value = {.simple_preference = REQUIRE}},
                [PRESERVE_MSG_BOUNDARIES] = {.value = {.simple_preference = REQUIRE}},
                [PER_MSG_RELIABILITY] = {.value = {.simple_preference = PREFER}},
                [PRESERVE_ORDER] = {.value = {.simple_preference = REQUIRE}},
                [ZERO_RTT_MSG] = {.value = {.simple_preference = NO_PREFERENCE}},
                [MULTISTREAMING] = {.value = {.simple_preference = NO_PREFERENCE}},
                [FULL_CHECKSUM_SEND] = {.value = {.simple_preference = REQUIRE}},
                [FULL_CHECKSUM_RECV] = {.value = {.simple_preference = REQUIRE}},
                [CONGESTION_CONTROL] = {.value = {.simple_preference = REQUIRE}},
                [KEEP_ALIVE] = {.value = {.simple_preference = NO_PREFERENCE}},
                [INTERFACE] = {.value = {.preference_set_val = {0}}},
                [PVD] = {.value = {.preference_set_val = {0}}},
                [USE_TEMPORARY_LOCAL_ADDRESS] = {.value = {.simple_preference = NO_PREFERENCE}},
                [MULTIPATH] = {.value = {.enum_val = CT_MULTIPATH_DISABLED}},
                [ADVERTISES_ALT_ADDRESS] = {.value = {.simple_preference = NO_PREFERENCE}},
                [DIRECTION] = {.value = {.enum_val = CT_DIRECTION_BIDIRECTIONAL}},
                [SOFT_ERROR_NOTIFY] = {.value = {.simple_preference = NO_PREFERENCE}},
            }
        },
        .init = quic_init,
        .send = quic_send,
        .init_with_send = quic_init_with_send,
        .listen = quic_listen,
        .close_listener = quic_close_listener,
        .close_connection = quic_close_connection,
        .close_socket = quic_close_socket,
        .abort = quic_abort,
        .clone_connection = quic_clone_connection,
        .free_socket_state = quic_free_socket_state,
        .close_connection_group = quic_close_connection_group,
        .set_connection_priority = quic_set_connection_priority,
        .free_connection_state = quic_free_state,
        .free_connection_group_state = ct_free_quic_connection_group_state
};
    \end{minted}
    \caption{QUIC protocol implementation struct}
    \label{lst:quic-impl-struct}
\end{listing}
% \lstinputlisting[language=C, caption={udp.c full implementation}]{src/udp.c}

\begin{listing}[htbp]
    \begin{minted}{C}
void on_socket_timer(uv_timer_t* timer_handle) {
    ct_quic_socket_state_t* socket_state = (ct_quic_socket_state_t*)timer_handle->data;
    if (!socket_state || !socket_state->picoquic_ctx) {
        log_error("QUIC context timer triggered but context is invalid");
        return;
    }

    log_trace("QUIC context timer triggered, checking for new QUIC packets to send");

    picoquic_quic_t* picoquic_ctx = socket_state->picoquic_ctx;
    size_t send_length = 0;

    struct sockaddr_storage from_address = {0};
    struct sockaddr_storage to_address = {0};
    int if_index = 0;
    picoquic_cnx_t* last_cnx = NULL;

    do {
        send_length = 0;

        // Allocate buffer on heap for each packet (freed in send callback)
        unsigned char* send_buffer_base = malloc(MAX_QUIC_PACKET_SIZE);
        // ... allocation error checking ...

        int rc =
            picoquic_prepare_next_packet(picoquic_ctx, picoquic_get_quic_time(picoquic_ctx),
                                         send_buffer_base, MAX_QUIC_PACKET_SIZE, &send_length,
                                         &to_address, &from_address, &if_index, NULL, &last_cnx);
        // ... error handling ...

        log_trace("Prepared QUIC packet of length %zu", send_length);
        if (send_length > 0) {
            int rc =
                ct_quic_send_packet_with_pktinfo(socket_state->poll_handle, send_buffer_base,
                                                 send_length, &from_address, &to_address, if_index);
            // ... send status handling and cleanup ...
        } else {
            log_trace("No QUIC data to send at this time");
            // No data to send, free the buffer
            free(send_buffer_base);
        }
    } while (send_length > 0);
    log_trace("Finished sending QUIC packets");

    reset_quic_timer(socket_state);
}
    \end{minted}
    \caption[libuv Timer Callback Function]{libuv timer callback function which checks in on Picoquic when required.}
    \label{lst:timer-callback-function}
\end{listing}

\chapter{Doxygen Documentation}
\label{app:doxygen}
% If you have exported Doxygen to LaTeX, you can input it here
% Or describe the API generated by Doxygen
The following documentation was generated via Doxygen for the CTaps internal API.
% \input{doxygen_output/refman.tex}
